\input{../../../templates/es/preambulo.tex}

% Los listados se toman de src/ con \lstinputlisting; la diapositiva es el archivo.
\lstset{
    basicstyle=\ttfamily\fontsize{8pt}{9.6pt}\selectfont,
    breaklines=true,
    breakatwhitespace=true,
    columns=fullflexible,
    keepspaces=true,
    xleftmargin=0.3em,
    framesep=2pt,
    aboveskip=0pt,
    belowskip=0pt
}

\newcommand{\marcoresultado}[3]{%
    \begin{frame}{#1}
        \begin{center}
            \includegraphics[height=0.62\textheight,width=\textwidth,keepaspectratio]{#2}
        \end{center}
        \vspace{0.25em}
        {\small #3\par}
    \end{frame}%
}

\newcommand{\marcoafirmacion}[2]{%
    \begin{frame}{#1}
        {\small #2\par}
    \end{frame}%
}

\date{}

\begin{document}

\tituloleccion{11}{Otros Problemas Comunes}

\section{Esta lección}

\begin{frame}{Esta lección}
La Lección 08 estableció que por lo general no se puede saber si una corrida tuvo éxito. Estos dos problemas son el caso raro en que una comprobación lo zanja.

\vspace{0.6em}
\begin{itemize}
    \item Un objetivo verificable puede ser una regla de parada (residual 0, conflictos 0).
    \item Cuando la restricción es el objetivo, el paisaje cerca de la respuesta es una meseta.
    \item Una caja finita de enteros se puede enumerar; eso responde preguntas que un AG no puede.
\end{itemize}
\end{frame}

% ================= coloreado_01_el_grafo =================
\begin{frame}{Testigos y certificados}
\[
R(x,y,z)=|f|+|g|+|w|,\quad (x,y,z)\in[-20,20]^3\cap\mathbb Z^3,
\quad |\mathcal X|=41^3=68{,}921,
\]
\[
C(c)=\sum_{(u,v)\in E}I[c_u=c_v],\quad
|\mathcal C|=3^{30}=205{,}891{,}132{,}094{,}649,\quad
\mathbb E[C]=|E|/3.
\]
Residuo o conflictos cero es un testigo. La enumeración y el retroceso exacto
añaden certificados de unicidad o imposibilidad que una corrida exitosa no da.
\end{frame}

\section{Cuando la restricción es el objetivo completo}

\begin{frame}{Cuando la restricción es el objetivo completo}
El coloreado de grafos es lo opuesto al sistema de ecuaciones. Allí el objetivo era
una cantidad más o menos suave a minimizar; aquí lo único que se mide es una
restricción violada. La aptitud es el número negado de aristas cuyos dos extremos comparten
un color, por lo que un coloreado legal y un coloreado óptimo son el mismo objeto y
no hay nada que optimizar una vez que lo alcanzas.
\end{frame}

\marcoresultado{Cuando la restricción es el objetivo completo}{../figures/coloreado_01_el_grafo.png}{%
    30 vértices, 64 aristas distintas, grado medio 4.27, y 3³⁰ = 205,891,132,094,649 coloreados. Un coloreado aleatorio en la semilla 7 rompe 26 de las 64 aristas}

% ================= coloreado_02_poblacion_aleatoria =================
\section{La línea base que provee la suerte}

\begin{frame}{La línea base que provee la suerte}
Un coloreado aleatorio asigna cada uno de los tres colores de forma independiente, así que cada arista
tiene una posibilidad entre tres de romperse. Eso da una cantidad esperada de conflictos
de |E|/3 antes de ejecutar nada, y la muestra debería acercarse a ello. Lo que importa
es el otro extremo del histograma: el mejor de 500 coloreados aleatorios, y qué
tan lejos está todavía del 0 que exige el problema.
\end{frame}

\begin{frame}[fragile]{(1) poblacion\_aleatoria()}
\lstinputlisting[basicstyle=\ttfamily\fontsize{8pt}{9.6pt}\selectfont,firstline=86,lastline=90]{../src/coloreado_02_poblacion_aleatoria.py}
\end{frame}

\begin{frame}[fragile]{(2) la muestra}
\lstinputlisting[basicstyle=\ttfamily\fontsize{8pt}{9.6pt}\selectfont,firstline=93,lastline=101]{../src/coloreado_02_poblacion_aleatoria.py}
\end{frame}

\marcoresultado{La línea base que provee la suerte}{../figures/coloreado_02_poblacion_aleatoria.png}{%
    /3 = 21.33`, y la mejor extracción aún rompe 11 aristas (17.2\%)}

% ================= coloreado_03_operadores_basados_en_aptitud =================
\section{Operadores que consultan la función de aptitud}

\begin{frame}{Operadores que consultan la función de aptitud}
En una meseta, un operador ciego es como lanzar una moneda. La respuesta de Gridin para este problema
es dejar que los operadores miren antes de saltar: una cruza que conserva los mejores
dos de \{ambos hijos, ambos padres\}, y una mutación que prueba algunos recoloreados
aleatorios y se queda con uno solo si reduce el conteo de conflictos. Ambos son
codiciosos (greedy), y la codicia es una elección real con un costo real, así que hay que medirlo en lugar de
solo afirmarlo.
\end{frame}

\begin{frame}[fragile]{(1) cruza\_n\_puntos()}
\lstinputlisting[basicstyle=\ttfamily\fontsize{8pt}{9.6pt}\selectfont,firstline=95,lastline=107]{../src/coloreado_03_operadores_basados_en_aptitud.py}
\end{frame}

\begin{frame}[fragile]{(2) cruza\_basada\_en\_aptitud\_n\_puntos()}
\lstinputlisting[basicstyle=\ttfamily\fontsize{8pt}{9.6pt}\selectfont,firstline=110,lastline=118]{../src/coloreado_03_operadores_basados_en_aptitud.py}
\end{frame}

\begin{frame}[fragile]{(3) mutacion\_basada\_en\_aptitud\_cambio\_aleatorio()}
\lstinputlisting[basicstyle=\ttfamily\fontsize{8pt}{9.6pt}\selectfont,firstline=121,lastline=133]{../src/coloreado_03_operadores_basados_en_aptitud.py}
\end{frame}

\begin{frame}[fragile]{(4) el experimento}
\lstinputlisting[basicstyle=\ttfamily\fontsize{7pt}{8.5pt}\selectfont,firstline=136,lastline=162]{../src/coloreado_03_operadores_basados_en_aptitud.py}
\end{frame}

\marcoresultado{Operadores que consultan la función de aptitud}{../figures/coloreado_03_operadores_basados_en_aptitud.png}{%
    La cruza sencilla de dos puntos mueve la media de la población en +0.03 aristas; la versión basada en aptitud en −2.32. El recoloreado codicioso mueve la media en −1.02 pero deja intactos 206 de los 500 individuos, y su mejor individuo (11) es *peor* que el mejor de un recoloreado ciego (9)}

% ================= coloreado_04_busqueda_completa =================
\section{La búsqueda que no termina, y la que sí}

\begin{frame}{La búsqueda que no termina, y la que sí}
Este es el paso para el que se construyó la lección. El algoritmo genético se ensambla y
se ejecuta, y en la mayoría de las semillas se detiene a una o dos aristas de un coloreado legal y
se queda ahí por el resto de su presupuesto - la meseta que el paso 1 predijo.
Luego, diez líneas de backtracking resuelven la misma pregunta de inmediato, y también prueban
que dos colores son imposibles, lo cual ninguna cantidad de ejecución del AG podría haber
establecido.
\end{frame}

\begin{frame}[fragile]{(1) seleccion\_rango\_con\_elite()}
\lstinputlisting[basicstyle=\ttfamily\fontsize{7pt}{8.5pt}\selectfont,firstline=139,lastline=158]{../src/coloreado_04_busqueda_completa.py}
\end{frame}

\begin{frame}[fragile]{(2) ejecutar()}
\lstinputlisting[basicstyle=\ttfamily\fontsize{7pt}{8.5pt}\selectfont,firstline=161,lastline=187]{../src/coloreado_04_busqueda_completa.py}
\end{frame}

\begin{frame}[fragile]{(3) coloreado\_exacto()}
\lstinputlisting[basicstyle=\ttfamily\fontsize{6pt}{7.2pt}\selectfont,firstline=190,lastline=227]{../src/coloreado_04_busqueda_completa.py}
\end{frame}

\begin{frame}[fragile]{(4) el barrido de semillas}
\lstinputlisting[basicstyle=\ttfamily\fontsize{8pt}{9.6pt}\selectfont,firstline=230,lastline=241]{../src/coloreado_04_busqueda_completa.py}
\end{frame}

\marcoresultado{La búsqueda que no termina, y la que sí}{../figures/coloreado_04_busqueda_completa.png}{%
    \textbf{1 de 4 semillas} alcanza un coloreado legal (semilla 11, generación 15). Las semillas 7 y 16 se estancan en 2 conflictos y la semilla 1 en 1 conflicto, todas aún atascadas después de las 200 generaciones completas. El backtracking encuentra un 3-coloreado adecuado en 40 nodos de búsqueda y prueba que no existe ningún 2-coloreado en 4 nodos — 0.69 ms para ambos}

% ================= ecuaciones_01_el_sistema =================
\section{Un problema cuya respuesta puede verificarse}



\marcoresultado{Un problema cuya respuesta puede verificarse}{../figures/ecuaciones_01_el_sistema.png}{%
    Los genes son enteros en \texttt{[-20, 20]}, así que el espacio completo es de 68,921 tripletas. Los residuales son enteros, no flotantes: a lo largo de la línea \texttt{(y, z) = (1, 1)} van de 3 dígitos en \texttt{x = 2} a 31 dígitos en \texttt{x = 20}}

% ================= ecuaciones_02_poblacion_aleatoria =================
\section{Lo que vale una muestra aleatoria aquí}

\begin{frame}{Lo que vale una muestra aleatoria aquí}
Antes de ejecutar un algoritmo genético, averigua qué se logra solo con suerte. En la lección 09
el muestreo aleatorio fue una línea base respetable. Aquí no lo es, y la razón es
visible en un histograma: el residual se mide en dígitos decimales, por lo que una extracción
que parece cercana sigue estando astronómicamente lejos.
\end{frame}

\begin{frame}[fragile]{(1) tripleta\_aleatoria()}
\lstinputlisting[basicstyle=\ttfamily\fontsize{8pt}{9.6pt}\selectfont,firstline=75,lastline=79]{../src/ecuaciones_02_poblacion_aleatoria.py}
\end{frame}

\begin{frame}[fragile]{(2) la muestra}
\lstinputlisting[basicstyle=\ttfamily\fontsize{8pt}{9.6pt}\selectfont,firstline=82,lastline=92]{../src/ecuaciones_02_poblacion_aleatoria.py}
\end{frame}

\marcoresultado{Lo que vale una muestra aleatoria aquí}{../figures/ecuaciones_02_poblacion_aleatoria.png}{%
    400 extracciones uniformes (0.58\% de la caja) encuentran 0 soluciones; el mejor es \texttt{(0, 1, 2)} con residual 453, mientras que la extracción mediana tiene un residual de 957 dígitos y la peor 14,187 dígitos}

% ================= ecuaciones_03_ag_completo =================
\section{El algoritmo genético lo resuelve}

\begin{frame}{El algoritmo genético lo resuelve}
Nada aquí es maquinaria nueva: la selección por rango con elitismo vino de la lección 03,
la cruza blend de la lección 04, la desviación gaussiana de la lección 05. Lo que es nuevo
es la regla de parada. Todas las demás lecciones se detuvieron después de un presupuesto fijo porque
no podían decir si habían ganado. Esta se detiene cuando el residual es cero,
porque eso es una prueba.
\end{frame}

\begin{frame}[fragile]{(1) Individuo}
\lstinputlisting[basicstyle=\ttfamily\fontsize{8pt}{9.6pt}\selectfont,firstline=88,lastline=99]{../src/ecuaciones_03_ag_completo.py}
\end{frame}

\begin{frame}[fragile]{(2) seleccion\_rango\_con\_elite()}
\lstinputlisting[basicstyle=\ttfamily\fontsize{7pt}{8.5pt}\selectfont,firstline=102,lastline=123]{../src/ecuaciones_03_ag_completo.py}
\end{frame}

\begin{frame}[fragile]{(3) cruza\_blend()}
\lstinputlisting[basicstyle=\ttfamily\fontsize{8pt}{9.6pt}\selectfont,firstline=126,lastline=139]{../src/ecuaciones_03_ag_completo.py}
\end{frame}

\begin{frame}[fragile]{(4) mutacion\_desviacion\_aleatoria()}
\lstinputlisting[basicstyle=\ttfamily\fontsize{8pt}{9.6pt}\selectfont,firstline=142,lastline=151]{../src/ecuaciones_03_ag_completo.py}
\end{frame}

\begin{frame}[fragile]{(5) ejecutar()}
\lstinputlisting[basicstyle=\ttfamily\fontsize{6pt}{7.2pt}\selectfont,firstline=154,lastline=187]{../src/ecuaciones_03_ag_completo.py}
\end{frame}

\marcoresultado{El algoritmo genético lo resuelve}{../figures/ecuaciones_03_ag_completo.png}{%
    La semilla 3 alcanza un residual de 0 en la generación 5 después de 2,794 evaluaciones, en \texttt{x = -6, y = 2, z = 3}; \texttt{f}, \texttt{g} y \texttt{w} son cada una individualmente 0}

% ================= ecuaciones_04_busqueda_exhaustiva =================
\section{El método que no necesita un algoritmo genético}

\begin{frame}{El método que no necesita un algoritmo genético}
Los genes son enteros y la caja es finita, así que todo el espacio de búsqueda puede
ser recorrido. Este es el paso donde el curso finalmente hace la pregunta que ha estado
evadiendo: dado que un algoritmo genético es una heurística, ¿cuándo es la herramienta
incorrecta? La respuesta aquí no es la que sugieren los conteos de evaluación, y el
paso está escrito alrededor de lo que el código realmente mide.
\end{frame}

\begin{frame}[fragile]{(1) enumerar\_caja()}
\lstinputlisting[basicstyle=\ttfamily\fontsize{7pt}{8.5pt}\selectfont,firstline=179,lastline=199]{../src/ecuaciones_04_busqueda_exhaustiva.py}
\end{frame}

\begin{frame}[fragile]{(2) ejecutar()}
\lstinputlisting[basicstyle=\ttfamily\fontsize{8pt}{9.6pt}\selectfont,firstline=145,lastline=145]{../src/ecuaciones_04_busqueda_exhaustiva.py}
\end{frame}

\begin{frame}[fragile]{(3) la comparación}
\lstinputlisting[basicstyle=\ttfamily\fontsize{8pt}{9.6pt}\selectfont,firstline=202,lastline=216]{../src/ecuaciones_04_busqueda_exhaustiva.py}
\end{frame}

\marcoresultado{El método que no necesita un algoritmo genético}{../figures/ecuaciones_04_busqueda_exhaustiva.png}{%
    Las 5 semillas probadas encuentran la misma raíz, en 5–28 generaciones y 2,794–13,764 evaluaciones. Enumerar la caja cuesta 68,921 evaluaciones (alrededor de 10× más, aproximadamente 3.7 s) y devuelve estrictamente más: hay \textbf{exactamente una} raíz en la caja}

\section{Siguiente}

\begin{frame}{Siguiente}
El sistema tiene exactamente una raíz en la caja. El AG la encuentra más barato que la enumeración, pero no prueba unicidad. Colorear tiene 205 billones de asignaciones, y el retroceso exacto certifica tres colores y descarta dos.

\vspace{0.8em}
\textbf{Lección 12 --- El algoritmo genético adaptativo}

Los ajustes fijos afinados de la 07, el TSP de la 10, presupuesto igual de evaluaciones. Adaptar es un seguro contra una constante mala, no un reemplazo del ajuste.
\end{frame}
\end{document}
